% ============================================================
%  TEMA 2 — APARTADO 2: DESCRIPCIÓN Y ESPECIFICACIÓN DE PUESTOS
%
%  RECORDATORIO AL EQUIPO:
%  Hay que presentar TRES puestos distintos de la empresa.
%  Para cada puesto debéis elaborar:
%
%  A) DESCRIPCIÓN DEL PUESTO (el QUÉ se hace):
%     - Identificación: denominación, departamento, nivel jerárquico, superior
%       inmediato, subordinados, fecha de elaboración.
%     - Resumen del puesto (misión en 2-3 líneas).
%     - Funciones y tareas principales (ordenadas por importancia o frecuencia).
%     - Responsabilidades: personas, equipos, presupuesto, información.
%     - Condiciones de trabajo: horario, lugar, movilidad, riesgos laborales.
%     - Relaciones internas y externas.
%
%  B) ESPECIFICACIÓN DEL PUESTO (el QUIÉN lo puede hacer):
%     - Formación mínima requerida.
%     - Experiencia mínima.
%     - Conocimientos específicos.
%     - Competencias / habilidades (técnicas y conductuales).
%     - Rasgos de personalidad relevantes.
%
%  PARA CADA PUESTO VALORAD:
%  - ¿La descripción está actualizada? ¿Refleja la realidad del puesto?
%  - ¿La especificación es adecuada (ni sobre- ni infracualificación)?
%  - ¿Está la empresa satisfecha con estas fichas?
%  - ¿Podría mejorarse? ¿Cómo?
% ============================================================

\section{Descripción y especificación de tres puestos de CaixaBank}

% Elegid tres puestos de diferentes niveles jerárquicos o áreas funcionales
% para dar variedad al análisis. Ejemplos posibles:
%   - Director/a de Oficina (nivel directivo / mando intermedio)
%   - Gestor/a de Banca Privada (nivel técnico especializado)
%   - Asesor/a Comercial de Oficina (nivel operativo)
% Ajustadlos a los puestos sobre los que hayáis conseguido información real.

% ===========================================================
\subsection{Puesto 1: [Nombre del puesto — p. ej. Director/a de Oficina]}
% ===========================================================

\subsubsection{Descripción del puesto}

\begin{table}[H]
  \centering
  \small
  \caption{Ficha de identificación — Puesto 1}
  \begin{tabularx}{\textwidth}{lX}
    \toprule
    \textbf{Campo}               & \textbf{Contenido} \\
    \midrule
    Denominación del puesto      & [Director/a de Oficina] \\
    Código / referencia interna  & [Si existe] \\
    Departamento / Área          & [Red Territorial / Oficinas] \\
    Nivel jerárquico             & [Mando intermedio] \\
    Superior inmediato           & [Director/a Territorial de Zona] \\
    Número de subordinados       & [$\sim$X empleados] \\
    Fecha de elaboración         & [Mes y año de la ficha original] \\
    Última revisión              & [Fecha] \\
    \bottomrule
  \end{tabularx}
\end{table}

\textbf{Misión del puesto:}

[Resumen en 2-3 líneas sobre la razón de ser del puesto en la empresa.
Ejemplo: ``Gestionar, coordinar y liderar el equipo de la oficina bancaria,
asegurando el cumplimiento de los objetivos comerciales y la satisfacción
del cliente, en el marco de las políticas y procedimientos del Grupo CaixaBank.'']

\textbf{Funciones y tareas principales:}
\begin{enumerate}
  \item [Tarea 1 — descripción breve]
  \item [Tarea 2]
  \item [Tarea 3]
  \item [...]
\end{enumerate}

\textbf{Responsabilidades:}
\begin{itemize}
  \item \textit{Sobre personas:} [supervisión, evaluación, formación del equipo]
  \item \textit{Sobre equipos / materiales:} [...]
  \item \textit{Sobre información confidencial:} [datos de clientes, riesgo crediticio]
  \item \textit{Sobre presupuesto:} [gestión de la P\&L de la oficina]
\end{itemize}

\textbf{Condiciones de trabajo:}
\begin{itemize}
  \item Horario: [jornada completa, partido / continua, ...]
  \item Lugar: [oficina física, posible teletrabajo parcial]
  \item Movilidad / desplazamientos: [sí/no; periodicidad]
  \item Riesgos laborales: [estrés, exposición prolongada a pantallas, ...]
\end{itemize}

\subsubsection{Especificación del puesto}

\begin{table}[H]
  \centering
  \small
  \caption{Requisitos del candidato — Puesto 1}
  \begin{tabularx}{\textwidth}{lX}
    \toprule
    \textbf{Requisito}       & \textbf{Descripción} \\
    \midrule
    Formación mínima         & [Ej. Grado en ADE, Economía, Derecho o similar] \\
    Formación complementaria & [Ej. Máster en Finanzas, certificaciones bancarias (EFA, EFP)] \\
    Experiencia mínima       & [Ej. 5 años en sector bancario, 2 en gestión de equipos] \\
    Conocimientos técnicos   & [Productos bancarios, normativa MiFID II, gestión de riesgos, ...] \\
    Idiomas                  & [Español nativo; inglés nivel B2 valorable] \\
    Competencias conductuales & [Liderazgo, orientación al cliente, toma de decisiones, ...] \\
    Herramientas / software  & [CRM bancario interno, Office 365, ...] \\
    \bottomrule
  \end{tabularx}
\end{table}

\begin{notaequipo}
  Valorad si la especificación de este puesto es realista o si exige
  requisitos excesivos (sobre-cualificación) o insuficientes. ¿Cambia mucho
  la especificación real de lo que publica CaixaBank en sus ofertas de
  empleo (InfoJobs, LinkedIn, web corporativa)?
\end{notaequipo}

% ===========================================================
\subsection{Puesto 2: [Nombre del puesto — p. ej. Gestor/a de Banca Privada]}
% ===========================================================

\subsubsection{Descripción del puesto}

\begin{table}[H]
  \centering
  \small
  \caption{Ficha de identificación — Puesto 2}
  \begin{tabularx}{\textwidth}{lX}
    \toprule
    \textbf{Campo}               & \textbf{Contenido} \\
    \midrule
    Denominación del puesto      & [Gestor/a de Banca Privada] \\
    Departamento / Área          & [Banca Privada / CaixaBank Private Banking] \\
    Nivel jerárquico             & [Técnico especializado] \\
    Superior inmediato           & [Director/a de Banca Privada] \\
    Número de subordinados       & [Ninguno / gestión de cartera propia] \\
    Fecha de elaboración / revisión & [Fecha] \\
    \bottomrule
  \end{tabularx}
\end{table}

\textbf{Misión del puesto:}

[Descripción a completar]

\textbf{Funciones y tareas principales:}
\begin{enumerate}
  \item [...]
\end{enumerate}

\textbf{Responsabilidades:}
\begin{itemize}
  \item [...]
\end{itemize}

\textbf{Condiciones de trabajo:}
\begin{itemize}
  \item [...]
\end{itemize}

\subsubsection{Especificación del puesto}

\begin{table}[H]
  \centering
  \small
  \caption{Requisitos del candidato — Puesto 2}
  \begin{tabularx}{\textwidth}{lX}
    \toprule
    \textbf{Requisito}       & \textbf{Descripción} \\
    \midrule
    Formación mínima         & [Ej. Grado en ADE / Economía] \\
    Certificaciones          & [Ej. CFA, EFP, EFPA nivel III] \\
    Experiencia mínima       & [Ej. 3-5 años en gestión de patrimonios] \\
    Conocimientos técnicos   & [Mercados financieros, fiscalidad de inversiones, ...] \\
    Competencias conductuales & [Empatía, discreción, orientación al resultado, ...] \\
    \bottomrule
  \end{tabularx}
\end{table}

\begin{notaequipo}
  Comparad este puesto con el anterior. ¿Hay diferencias claras en el nivel
  de exigencia? ¿La empresa diferencia bien entre banca personal y banca privada?
\end{notaequipo}

% ===========================================================
\subsection{Puesto 3: [Nombre del puesto — p. ej. Asesor/a Comercial de Oficina]}
% ===========================================================

\subsubsection{Descripción del puesto}

\begin{table}[H]
  \centering
  \small
  \caption{Ficha de identificación — Puesto 3}
  \begin{tabularx}{\textwidth}{lX}
    \toprule
    \textbf{Campo}               & \textbf{Contenido} \\
    \midrule
    Denominación del puesto      & [Asesor/a Comercial de Oficina] \\
    Departamento / Área          & [Red Comercial] \\
    Nivel jerárquico             & [Operativo / técnico base] \\
    Superior inmediato           & [Director/a de Oficina] \\
    Número de subordinados       & [Ninguno] \\
    Fecha de elaboración / revisión & [Fecha] \\
    \bottomrule
  \end{tabularx}
\end{table}

\textbf{Misión del puesto:}

[Descripción a completar]

\textbf{Funciones y tareas principales:}
\begin{enumerate}
  \item [...]
\end{enumerate}

\textbf{Responsabilidades:}
\begin{itemize}
  \item [...]
\end{itemize}

\textbf{Condiciones de trabajo:}
\begin{itemize}
  \item [...]
\end{itemize}

\subsubsection{Especificación del puesto}

\begin{table}[H]
  \centering
  \small
  \caption{Requisitos del candidato — Puesto 3}
  \begin{tabularx}{\textwidth}{lX}
    \toprule
    \textbf{Requisito}       & \textbf{Descripción} \\
    \midrule
    Formación mínima         & [Ej. FP Superior en Administración, Grado ADE o similar] \\
    Certificaciones          & [EFA, EFB o formación interna CaixaBank] \\
    Experiencia mínima       & [0-2 años; puede ser puesto de entrada] \\
    Conocimientos técnicos   & [Productos bancarios básicos, normativa, CRM, ...] \\
    Competencias conductuales & [Orientación al cliente, proactividad, comunicación, ...] \\
    \bottomrule
  \end{tabularx}
\end{table}

\begin{notaequipo}
  Este suele ser el puesto de entrada más común. Analizad si la rotación
  es alta en este nivel y si eso refleja problemas en la especificación
  (se selecciona a perfiles que luego no encajan) o en las condiciones del puesto.
  ¿Cuánto tarda CaixaBank en cubrir una vacante de este tipo?
\end{notaequipo}
